% !TEX root = ../Thesis.tex

\chapter{基于公开数据集的初步研究}\label{chap:preliminary}

{
本章介绍利用公开的光学散斑数据集进行的前期探索工作，通过复现已有研究结果，验证深度学习方法在多模光纤散斑图像重建中的可行性。

\section{引言}

% TODO: 简要说明本章的研究目的和意义


\section{公开数据集介绍}

\subsection{数据集来源}

% TODO: 介绍使用的公开数据集，包括数据来源、文献引用等

\subsection{数据集特征}

% TODO: 描述数据集的规模、图像尺寸、散斑特征等

\subsection{数据预处理}

% TODO: 介绍对数据集进行的预处理操作


\section{网络模型设计与训练}

\subsection{网络架构选择}

% TODO: 说明选择UNet网络的原因

\subsection{网络参数设置}

% TODO: 详细列出网络的各项参数配置

\subsection{训练策略}

\subsubsection{损失函数}

% TODO: 介绍使用的损失函数

\subsubsection{优化器}

% TODO: 介绍优化器的选择和参数设置

\subsubsection{训练过程}

% TODO: 描述训练的具体过程


\section{实验结果与分析}

\subsection{重建结果展示}

% TODO: 展示图像重建的结果（图片）

\subsection{定量评价}

% TODO: 使用PSNR、SSIM等指标进行定量评价

\subsection{与原文献结果的对比}

% TODO: 将复现结果与原论文进行对比分析


\section{影响因素分析}

\subsection{网络深度的影响}

% TODO: 分析网络深度对重建效果的影响

\subsection{训练数据量的影响}

% TODO: 分析训练数据量对性能的影响

\subsection{其他参数的影响}

% TODO: 分析其他关键参数的影响


\section{初步结论}

% TODO: 总结本章的探索工作得到的结论，为后续研究提供基础


\section{本章小结}

% TODO: 总结本章内容

}

